% GENERATED FILE. Edit canonical lessons, then run npm run generate:books.
% canonical-source-hash: fnv1a64:734f67e4b91f1104
% canonical-lessons: PA-C198-message-return, PA-C198-suneha, PA-C198-suneha-review, PA-C198-message-label, PA-C198-payoff

\chapter{This Is a Message}
\label{ch:pa-a2-suneha}
\hlchaptermodality{198}

\begin{chapteropening}
\textbf{By the end of this chapter:} \emph{I can recognise and use the noun for a message in one bounded identity frame.}

Recognise and produce the message noun in one bounded identity frame through listening, speaking, reading, and writing.
\end{chapteropening}

\section[\texorpdfstring{\pa{ਇਹ} · \pa{ਹੈ} (eh; hai)}{ਇਹ · ਹੈ (eh; hai)}]{Return to this and is}

\label{lesson:PA-C198-message-return}



\begin{quote}
\emph{Hear:} \emph{eh; hai}; YOU SAY: \textbf{\pa{ਇਹ}; \pa{ਹੈ}}
\end{quote}

\begin{grammarlens}[title={Grammar Lens}]
Keep the familiar pointer first and the familiar verb last. The middle slot stays empty until the next lesson owns its new word.
\end{grammarlens}

\subsection*{Your turn}
\begin{itemize}
\raggedright
  \item \emph{Hear:} \emph{eh}; YOU SAY: \textbf{\pa{ਇਹ}}
  \item \emph{Hear:} \emph{hai}; YOU SAY: \textbf{\pa{ਹੈ}}
  \item Say the first word, pause for the empty middle, then say the last.
\end{itemize}

\begin{tcolorbox}[breakable,colback=teal!4,colframe=teal!35!black,title={Before you move on}]
Produce the two familiar words around the empty middle slot. This opener owns no new word.
\end{tcolorbox}
\section[\texorpdfstring{\pa{ਸੁਨੇਹਾ} (sunehaa)}{ਸੁਨੇਹਾ (sunehaa)}]{One noun for a message}

\label{lesson:PA-C198-suneha}



\begin{quote}
Retrieve the familiar first and last words of the identity frame.
\end{quote}

\begin{grammarlens}[title={Grammar Lens: direct owner}]
\begin{quote}
\textbf{\pa{ਸੁਨੇਹਾ}} — \emph{sunehā} — message; information communicated from one party to another
\end{quote}

Hear and say the whole word before reading it. Keep it as one whole noun rather than inventing smaller Punjabi pieces. This lesson owns the word alone; the complete identity frame comes only after a separate retrieval.
\end{grammarlens}

\begin{cousinweb}
Punjabi University's Punjabipedia defines this exact word as information given from one party to another. It supports only the present-day message sense taught here. Source: \href{https://punjabipedia.org/topic.aspx?txt=\%E0\%A8\%B8\%E0\%A9\%81\%E0\%A8\%A8\%E0\%A9\%87\%E0\%A8\%B9\%E0\%A8\%BE}{Punjabi University Punjabipedia}.

The exact history of modern written \textbf{\pa{ਸੁਨੇਹਾ}} remains unresolved in the sources used here. For a historical comparison only, Turner's CDIAL entry 12904 groups Punjabi \emph{sanehā̃} with Sanskrit \emph{saṁdēśa}, “message, information.” The source does not document the intermediate changes into this modern spelling. This historical comparison is not a letter-by-letter derivation and does not create a productive Punjabi rule: keep the present-day Punjabi noun whole. Source: \href{https://dsal.uchicago.edu/cgi-bin/app/soas\_query.py?qs=sa\%E1\%B9\%81d\%C4\%93\%C5\%9Ba\&searchhws=yes\&matchtype=exact}{Turner, CDIAL 12904}.
\end{cousinweb}

\subsection*{Your turn}
\begin{itemize}
\raggedright
  \item \emph{Hear:} \emph{sunehā}; YOU SAY: \textbf{\pa{ਸੁਨੇਹਾ}}
  \item \emph{Read it:} \textbf{\pa{ਸੁਨੇਹਾ}}
  \item \emph{Write it:} copy \textbf{\pa{ਸੁਨੇਹਾ}} once, checking the whole word
\end{itemize}

\begin{tcolorbox}[breakable,colback=teal!4,colframe=teal!35!black,title={Before you move on}]
Produce the noun for “message” without looking.
\end{tcolorbox}
\section[\texorpdfstring{\pa{ਸੁਨੇਹਾ} (sunehaa)}{ਸੁਨੇਹਾ (sunehaa)}]{Retrieve the new word on its own}

\label{lesson:PA-C198-suneha-review}



\begin{quote}
\emph{Hear:} \emph{sunehā}; YOU SAY: \textbf{\pa{ਸੁਨੇਹਾ}}
\end{quote}

\subsection*{Your turn}
\begin{itemize}
\raggedright
  \item \emph{Hear:} \emph{message}; YOU SAY: \textbf{\pa{ਸੁਨੇਹਾ}}
  \item \emph{Read it:} \textbf{\pa{ਸੁਨੇਹਾ}}
  \item \emph{Write it:} cover and write \textbf{\pa{ਸੁਨੇਹਾ}}
\end{itemize}

\begin{grammarlens}[title={Grammar Lens}]
Keep one noun and one message meaning here. Inflection, specialised digital uses, and a send-message construction stay outside this bounded lesson.
\end{grammarlens}

\begin{tcolorbox}[breakable,colback=teal!4,colframe=teal!35!black,title={Before you move on}]
Say and then write the word once from memory.
\end{tcolorbox}
\section[\texorpdfstring{\pa{ਇਹ} \pa{ਸੁਨੇਹਾ} \pa{ਹੈ} (eh sunehaa hai)}{ਇਹ ਸੁਨੇਹਾ ਹੈ (eh sunehaa hai)}]{Put the new noun into one identity frame}

\label{lesson:PA-C198-message-label}



\begin{quote}
Retrieve \textbf{\pa{ਇਹ}}, then the new noun, then \textbf{\pa{ਹੈ}}.
\end{quote}

\begin{grammarlens}[title={Grammar Lens}]
\begin{quote}
\textbf{\pa{ਇਹ} \pa{ਸੁਨੇਹਾ} \pa{ਹੈ}} — \emph{eh sunehā hai} — this is a message
\end{quote}

Keep the new whole noun between familiar \textbf{\pa{ਇਹ}} and \textbf{\pa{ਹੈ}}. Familiar \textbf{\pa{ਹੈ}} stays last in the already-owned identity frame. This chapter teaches only this one label; it does not add an inflection rule or a send-message construction.
\end{grammarlens}

\subsection*{Your turn}
\begin{itemize}
\raggedright
  \item \emph{Hear:} \emph{eh sunehā hai}; YOU SAY: \textbf{\pa{ਇਹ} \pa{ਸੁਨੇਹਾ} \pa{ਹੈ}}
  \item \emph{Read it:} \textbf{\pa{ਇਹ} \pa{ਸੁਨੇਹਾ} \pa{ਹੈ}}
  \item \emph{Write it:} cover and write \textbf{\pa{ਇਹ} \pa{ਸੁਨੇਹਾ} \pa{ਹੈ}}
\end{itemize}

\begin{tcolorbox}[breakable,colback=teal!4,colframe=teal!35!black,title={Before you move on}]
Produce the complete identity frame.
\end{tcolorbox}
\section[\texorpdfstring{\pa{ਇਹ} \pa{ਸੁਨੇਹਾ} \pa{ਹੈ} (eh sunehaa hai)}{ਇਹ ਸੁਨੇਹਾ ਹੈ (eh sunehaa hai)}]{A bounded message-noun payoff}

\label{lesson:PA-C198-payoff}



\begin{quote}
Retrieve the new noun alone before rebuilding its identity frame.
\end{quote}

\subsection*{Your turn: four skills}
\begin{itemize}
\raggedright
  \item \emph{Hear:} \emph{eh sunehā hai}; YOU SAY: \textbf{\pa{ਇਹ} \pa{ਸੁਨੇਹਾ} \pa{ਹੈ}}
  \item \emph{Read it:} \textbf{\pa{ਇਹ} \pa{ਸੁਨੇਹਾ} \pa{ਹੈ}}
  \item \emph{Write it:} cover the line and write \textbf{\pa{ਇਹ} \pa{ਸੁਨੇਹਾ} \pa{ਹੈ}}
\end{itemize}

\begin{tcolorbox}[breakable,colback=teal!4,colframe=teal!35!black,title={Before you move on}]
Say the new noun, then produce the complete identity frame.

This is the twenty-first prerequisite tranche exposed by Punjabi A2 mock drafting. It does not unblock Punjabi A2 mock 1 by itself. It is not an A2 mock, not A2 readiness or attainment, and not external or human validation.
\end{tcolorbox}
